%【说明】第二章按照方法类别组织相关工作，每一节均以与本文直接相关的研究缺口收束。
\section{Literature Review}
\label{sec:literature}

\subsection{Dynamic Scheduling in Manufacturing Systems}
\label{subsec:traditional_methods}

Dynamic scheduling methods differ mainly in when they react to an event and how they
construct the revised schedule. Reactive dispatching rules select an operation or resource
from the currently feasible set and therefore provide short decision times. Predictive-
reactive methods start from a baseline schedule and repair or regenerate it after a
disturbance. Mathematical programming and metaheuristic methods can explore a wider solution
space, whereas rolling-horizon and event-driven mechanisms restrict each optimization to the
information available at the current time.

%【TODO】为派工规则、预测反应式调度、滚动时域、数学规划和元启发式调度分别补充代表性原始文献。

These paradigms provide different balances between responsiveness and search quality.
Dispatching rules are inexpensive but encode only the information included in their priority
indices. Reoptimization can account for more constraints but may require substantial
computation when disturbances occur frequently. Population-based methods can produce diverse
solutions but repeatedly evaluating candidate schedules can limit their use at dense decision
epochs. The choice of method is therefore inseparable from the state representation and the
time available for online decisions.

Multi-level manufacturing-assembly systems introduce an additional coordination requirement.
A manufacturing delay changes not only the completion time of one route but also the release
time of an assembly that waits for several branches. Shared resources couple otherwise
separate routes, and successive kitting points propagate the effect toward higher assembly
levels. Existing dynamic scheduling formulations do not automatically preserve these paths
when the system is summarized by global workload and utilization indicators. This motivates
a relational representation that remains synchronized with the event-driven environment.

%【TODO】补充多级装配、齐套调度和制造装配协同调度的已核验文献，并明确本文问题与柔性作业车间调度的区别。

\subsection{Deep Reinforcement Learning for Dynamic Scheduling}
\label{subsec:drl_scheduling}

Deep reinforcement learning formulates scheduling as sequential interaction among a state,
an action, a transition mechanism, and a reward. Value-based methods estimate the long-term
return of discrete decisions, while policy-gradient and actor-critic methods optimize a
parameterized policy directly. Once trained, these approaches can produce decisions through
network inference and are therefore attractive for event-driven environments that require
repeated online responses.

%【TODO】补充价值型、策略梯度型和Actor-Critic调度方法的代表性原始文献。

The performance of a learned scheduler depends on how the Markov decision process is
constructed. State variables must expose information needed to distinguish feasible and
consequential actions. The action space must correspond to decisions that the simulator can
execute, and feasibility masking must prevent violations of precedence, kitting, and resource
capacity. The reward must remain aligned with the reported scheduling objectives. A policy
can otherwise converge in its training environment while optimizing a surrogate that does
not correspond to the final evaluation.

Generalization is another central issue. A policy trained on a limited set of instances may
associate decisions with a fixed number or ordering of operations. Relational encoders,
parameter sharing, randomized training scenarios, and independent test instances have been
used to reduce this dependence. Nevertheless, generalization claims require an explicit
train-test boundary and a precise description of what changes between the two partitions.
This requirement is particularly important when topology, scale, and disturbance intensity
can vary independently.

%【TODO】补充调度策略泛化、可行性掩码、奖励设计和独立测试协议方面的已核验文献。

\subsection{Graph Representation Learning for Scheduling}
\label{subsec:graph_scheduling}

Graph-based scheduling methods encode operations and resources as nodes and express
precedence or eligibility through edges. Message passing aggregates information from local
neighborhoods, which supports variable-size instances and preserves structural constraints
more directly than a fixed-length state vector. Attention mechanisms can further assign
different weights to neighboring entities when their current influence on a decision differs.

Homogeneous graphs use a shared transformation for all nodes or edges. This representation is
effective when the encoded relations have similar meanings, but manufacturing-assembly
scheduling contains several distinct semantics. A technological edge indicates an execution
order, an eligibility edge indicates a possible resource assignment, an occupation edge
records a current physical state, and a kitting edge synchronizes multiple branches. Treating
these relations identically can obscure why information from a neighbor should influence the
current operation.

Heterogeneous graph neural networks address this issue by using type-specific projections and
relation-specific message functions. For dynamic scheduling, the graph must additionally
distinguish persistent production structure from event-dependent features. An effective model
should update the affected portion of the graph after an event, retain the semantic identity
of each relation, and provide a graph-level representation for policy learning. These needs
motivate the dynamic heterogeneous representation and multi-relation semantic encoder used in
this study.

%【TODO】补充图调度、图注意力、异构图神经网络及动态图学习的代表性原始文献。

\subsection{Research Gaps and Motivation}
\label{subsec:graph_drl}

The literature indicates three connected gaps. First, common dynamic scheduling states do not
explicitly preserve the propagation of a local event through shared resources and successive
kitting levels. Second, graph-based schedulers often emphasize technological precedence and
resource eligibility but do not distinguish all relations required by multi-level assembly.
Third, experimental generalization is sometimes inferred from training performance without a
strictly separated test partition or a clear definition of the distribution shift.

The present study addresses these gaps through a common environment-representation-policy
interface. The discrete-event simulator defines how disturbances change physical feasibility.
The dynamic heterogeneous graph records both persistent structure and current event effects.
The multi-relation HGNN preserves semantic differences during message propagation, and PPO
optimizes the policy over the resulting state. The experiment design separates 30 training
instances from 15 independent test instances so that training behavior and transfer behavior
can be discussed independently.

%【TODO】完成文献检索后增加相关工作对比表，至少比较问题类型、扰动类型、装配关系、图类型、决策方法和泛化设置。
%【TODO】当前参考文献库仅含PPO条目，投稿前必须补齐本章所有文献并逐条核验DOI、作者、年份和方法归属。
